% ====================================================================
% s33_blend.tex — §3.3 혼합 음극: 공통-μ 대정준 + GS-2 비가산 공백 (W2 초안 — 실측·데이터 강조축)
%   중심식 = eq:sm-mc-balance 의 클래스곱→host 곱 한 줄 일반화(재유도 말고 xr). f_Si=Q_Si/Q∈[0,0.3].
%   강조: GS-2 비가산 실측 근거(BLEND_UP Ai2022·Chatzogiannakis2025 정밀 반영) + 블렌드 증거표.
% ====================================================================
\section{혼합 음극 --- 공통-$\mu$ 대정준과 비가산 실측 공백}\label{sec:si-blend}
흑연 다수에 Si 소수를 섞은 블렌드 음극은 두 활물질이 \emph{하나의} 집전체로 전자 회로를 공유하고 \emph{같은}
전해질에 잠긴다 --- 곧 두 host 는 같은 전기화학 저장조와 입자를 교환하며, 전극에 거는 전위 손잡이는
$U_\oc$ 하나뿐이다. 이 실험 구도가 그대로 통계역학 언어로 옮겨지는 것이 본 절의 중심이다: 다클래스 전하
보존 반전(\S\ref{sec:sm-mc} 식~\eqref{eq:sm-mc-balance})의 자리 클래스 집합을 \emph{두 host 클래스의 서로소
합집합}으로 넓히면, 블렌드 합성의 중심식이 한 줄로 나온다.

\subsection{공통-$\mu$ 대정준 반전 --- 클래스곱에서 host 곱으로}\label{ssec:blend-derivation}
\S\ref{sec:sm-mc} 는 한 host 의 자리를 전이 $j$ 별 클래스로 갈라 대정준 합을 클래스곱으로 인수분해했다
(식~\eqref{eq:sm-mc-factor}). 블렌드에서 자리 전체는 흑연 클래스와 Si 클래스의 서로소 합집합이고, 두
host 가 \emph{같은} 저장조(하나의 $\mu$)와 교환하므로 인수분해가 host 위에서 한 겹 더 갈라질 뿐이다:
\begin{equation}
\Xi=\prod_{\mathrm{host}\in\{\mathrm{gr},\mathrm{Si}\}}\ \prod_{j}\bigl(\Xi_{1,j}^{\mathrm{host}}\bigr)^{M_j^{\mathrm{host}}},
\qquad
\ln\Xi=\sum_{\mathrm{host}}\sum_{j}M_j^{\mathrm{host}}\ln\Xi_{1,j}^{\mathrm{host}}
\label{eq:blend-factor}
\end{equation}
--- 식~\eqref{eq:sm-mc-factor} 의 단일 index 집합을 host 로 라벨된 합집합으로 바꾼 것 그대로다(독립 부분계
분배함수 인수분해). 평균 입자수는 항별로 갈라져 host$\cdot$클래스마다 단일 자리 계산이 반복되고
(식~\eqref{eq:sm-mc-occ} 의 host 판), 자리당 전하 $F/N_A$ 로 읽어 $Q_j^{\mathrm{host}}\equiv(F/N_A)M_j^{\mathrm{host}}$,
$Q_\mathrm{gr}=\sum_jQ_j^{\mathrm{gr}}$, $Q_\mathrm{Si}=\sum_jQ_j^{\mathrm{Si}}$, $Q=Q_\mathrm{gr}+Q_\mathrm{Si}$
로 두면 고정 추출 전하 $\bar x$ 에서 관측 평형 전위 $U_\oc$ 를 정하는 보존식이 나온다:
\begin{equation}
\boxed{\;\sum_{\mathrm{host}\in\{\mathrm{gr},\mathrm{Si}\}}\ \sum_{j}Q_j^{\mathrm{host}}\,\xi_{\eq,j}^{\mathrm{host}}\bigl(U_\oc,T\bigr)\;=\;Q\,\bar x,
\qquad f_\mathrm{Si}=\frac{Q_\mathrm{Si}}{Q}\in[0,0.3]\;.}
\label{eq:blend-balance}
\end{equation}
이것이 블렌드의 전하 보존 중심식이며, 식~\eqref{eq:sm-mc-balance} 의 \emph{한 줄 일반화}다 --- 새 유도가
아니라 클래스곱을 host 곱으로 넓힌 것뿐이다. 반전의 지위(고정-함량 반전의 \emph{정의상 implicit
formulation}, 대정준$\to$정준 Legendre 켤레)와 유일근도 그대로 이월된다: 좌변은 강단조 logistic 의
용량가중 이중합이라 여전히 $U_\oc$ 에 순증하므로, 요동 양성 유일근 증명(식~\eqref{eq:sm-mc-fluc})이
문자 그대로 성립한다(host 가 둘이어도 합이 늘 뿐 단조성은 불변 --- 재증명 불필요).
\begin{verifybox}
회수 검산 --- $f_\mathrm{Si}=0$($Q_\mathrm{Si}{=}0$)이면 Si 합이 소멸해
식~\eqref{eq:blend-balance} 는 흑연 단독 보존식~\eqref{eq:sm-mc-balance}($Q{=}Q_\mathrm{gr}$)로 \emph{정확히}
되돌아간다 --- \S\ref{sec:si-code} 의 bit-exact 게이트가 이 회수를 코드에서 강제한다. 반대로 $f_\mathrm{Si}$
를 올리면 Si 클래스 용량 $Q_\mathrm{Si}=f_\mathrm{Si}Q$ 가 흑연 몫 $Q_\mathrm{gr}=(1-f_\mathrm{Si})Q$ 에서
넘어올 뿐, 좌변의 함수 구조는 불변이다.
\end{verifybox}

\subsection{블렌드 dQ/dV 합성 --- 공유 $U_\oc$ 위의 용량가중 합}\label{ssec:blend-dqdv}
식~\eqref{eq:blend-balance} 를 전위로 미분하면(서로소 자리 분해 --- 각 전이 진행이 독립), 관측 ICA 는
배경 위 두 host 기여의 선형 합으로 닫힌다:
\begin{equation}
\frac{\dd Q}{\dd V}(U_\oc)\;=\;C_\bg\;+\;\sum_{\mathrm{host}\in\{\mathrm{gr},\mathrm{Si}\}}\ \sum_{j}Q_j^{\mathrm{host}}\,
\frac{\dd\xi_{\eq,j}^{\mathrm{host}}}{\dd V}
\label{eq:blend-dqdv}
\end{equation}
--- 흑연 합산식~\eqref{eq:sum} 의 전이 합에 Si (유효) 전이 합이 \emph{같은 $U_\oc$ 축 위에} 더해진 것이다.
$f_\mathrm{Si}$ 는 두 host 봉우리 무리의 상대 면적을 가르는 손잡이이고($Q_\mathrm{Si}$ vs $Q_\mathrm{gr}$),
Si 봉우리들은 \S\ref{sec:si-survival} N5 의 재해석대로 \emph{유효 곡선 기술 성분}(진짜 staging 아님)이다.
저-Si 복합의 dQ/dV 에서 흑연 피크($\sim$0.2$\cdot$0.1$\cdot$0.07 V)와 Si 피크가 분리 관측된다는
실측\cite{naboka_sic2021} 이 이 합성의 전위축 접점을 준다. 이 식이 \S\ref{sec:si-code} 코드 합성 규칙의
수식 원형이다.

\subsection{정직 공백 GS-2 --- 비가산$\cdot$구간교대의 실측 편차}\label{ssec:blend-gs2}
식~\eqref{eq:blend-balance}$\cdot$\eqref{eq:blend-dqdv} 는 \emph{평형}($|I|\to0$)에서 정확한 전하 회계다 ---
두 host 가 한 전위를 공유한다는 것 외에 새 가정이 없다. 그러나 이를 \emph{유한율} 블렌드 dQ/dV 의 합성
규칙으로 쓰는 순간 두 암묵 가정이 얹힌다: \textbf{(A) 동시 co-reaction} --- 두 host 가 공유 $U_\oc$ 에서
동시에 반응한다; \textbf{(B) 가산성} --- 총 곡선이 두 단일-host 기여의 용량가중 단순합이다. 블렌드 실측은
이 둘에서 정성적으로 벗어난다.
\begin{enumerate}[label=\arabic*.,leftmargin=2.0em,itemsep=1pt]
\item \emph{구간별 host 교대}(가정 A 반례). 연속체 모델은 높은 SoC 에서 흑연이 반응전류의 대부분을
담당하다가 깊은 방전(low SoC)으로 갈수록 급격히 Si 상으로 전환됨을 보인다 --- 서로 다른 OCV
곡선$\cdot$질량분율$\cdot$교환전류밀도의 결과다\cite{ai_composite2022}. 곧 유한율에서 우세 host 는
\emph{구간별로 교대}하며 완전 동시반응이 아니다.
\item \emph{비가산 결합}(가정 B 반례). ``decoupled blend cell'' 로 성분을 분리 측정하면 혼합 전극 거동이
두 성분의 단순합이 \emph{아님}(비가산적)이 확인되고, 특히 탈리튬화 시 Si$\cdot$흑연 간 유효 C-rate 격차가
크다\cite{chatzogiannakis_blend2025} --- 공통-μ ``이상적 동시반응''과의 정량 편차다.
\item \emph{지지와 대가}(가정의 물리적 근거는 실재). 두 host 가 공유하는 전위 창은 실측된다: SiO$_x$/흑연에서
``겹치는 리튬화 전위(overlapping lithiation potentials)''가 직접 확인되며\cite{zhan_siox2026}, 이것이
공통-μ 의 물리적 근거다. 단 그 대가로 계면 균열과 히스 에너지손실이 SiO$_x$ 분율에 거의 \emph{선형}으로
증가한다\cite{zhan_siox2026}. 용량손실은 Si 함량과 상관되나 두께변화율과는 뚜렷한 상관이 없어(구분되는 두
열화축)\cite{moyassari_blend2022}, 블렌드가 단일 모드의 단순 중첩이 아님을 방증한다.
\end{enumerate}
곧 공통-μ 완전 동시반응$\cdot$가산 합성은 \emph{1차 근사}이며, 실측과 정성적 편차가 있다. 공백 분류(본문
공백 관행): \emph{물리 가정 충돌} --- 평형 가산-동시반응 가정이 실측의 비가산$\cdot$구간교대 동역학과
충돌한다. 논리 공백(식은 논리적으로 무결)도, 단순 수치 미해(평형 반전은 유일근 존재)도 아니다. 구간별 host
전환의 kinetic 결합 모델링은 \emph{범위 밖}으로 선언한다(GS-2 공백 명시까지 --- 구현 확장은 후속 버전).

\begin{table}[h]
\centering\footnotesize
\renewcommand{\arraystretch}{1.3}
\caption{흑연/Si 블렌드 실측$\cdot$모델 증거(BLEND\_UP 8건) --- 공통-μ 가정에 대한 함의. tier A $=$ 1차
정량 독립 확인$\cdot$B $=$ 서지 확정, 정량 초록 elided(확인 필요). f\_Si 스윕은 문헌 3$\cdot$4$\cdot$5 를
합쳐 $0$--$30$ wt\% 전 구간이 실측 커버된다.}
\label{tab:blend-evidence}
\begin{tabular}{@{}p{0.235\textwidth} p{0.135\textwidth} p{0.245\textwidth} p{0.235\textwidth}@{}}
\toprule
문헌 (tier) & f\_Si / 대상 & 측정$\cdot$모델 & 공통-μ 함의 \\
\midrule
tu\_blend2024 (A)            & SiO/흑연  & MSMR ROM(섭동, Verbrugge 계보) & 이론 다리(host별 MSMR) \\
ai\_composite2022 (A)        & Si/흑연   & 연속체 반응전류 모델            & 한계: 구간별 host 교대 \\
gautam\_blend2024 (A)        & 5--20 wt\%& 용량$\cdot$ICE 스윕            & f\_Si 초기값(Si15 최적) \\
moyassari\_blend2022 (A)     & 0--20 wt\%& 딜라토메트리$\cdot$용량손실     & 두 열화축 분리(비단순중첩) \\
chatzogiannakis\_blend2025 (A)& 30 wt\%   & decoupled blend cell          & 한계: 비가산$\cdot$C-rate 격차 \\
zhan\_siox2026 (A)           & SiO$_x$/흑연& 딜라토메트리$\cdot$DV/IC       & 지지+대가: 전위중첩$\cdot$균열 \\
zhang\_150ah2024 (B)         & SiO/흑연  & 150 Ah 상용 각형 셀           & 상용 스케일(정량 확인 필요) \\
mertin\_entropy2023 (B)      & Si/흑연   & 엔트로피 성분 해석             & 성분분리(정량 확인 필요) \\
\bottomrule
\end{tabular}
\end{table}

\begin{keybox}
\textbf{블렌드 한 줄.} 공통-μ 전하 보존 반전~\eqref{eq:blend-balance} 은 골격 이월(N9)이자 코드 합성
규칙이되 --- \emph{평형}에서 정확하고 $f_\mathrm{Si}{=}0$ 에서 흑연 단독으로 회수된다. 유한율 블렌드
dQ/dV 의 가산 합성은 \emph{1차 근사}이며(GS-2 --- 물리 가정 충돌), 실측이 구간별 host
교대\cite{ai_composite2022}$\cdot$비가산 결합\cite{chatzogiannakis_blend2025}을 보인다. 그 정밀 모델링은
범위 밖으로 정직히 선언한다.
\end{keybox}

% 자산(v1.0.22 R5 신규): [V22-R5-11 공통-μ 대정준 반전 host 곱 일반화 eq:blend-balance(=eq:sm-mc-balance 한 줄 확장·유일근 xr 이월)]
%   [V22-R5-12 f_Si=0 회수 검산·블렌드 dQ/dV 합성 eq:blend-factor·eq:blend-dqdv]
%   [V22-R5-13 GS-2 비가산·구간교대 정직 공백 ssec:blend-gs2(물리 가정 충돌·범위 밖 선언)]
%   [V22-R5-14 블렌드 실측 증거표 tab:blend-evidence(BLEND_UP 8건·공통-μ 함의·f_Si 0~30% 커버)]
